\documentclass[11pt]{article}
\usepackage{amsmath,amssymb}
\title{A raw tree quartic is not controlled by physical curvature energy}\author{}\date{}
\begin{document}\maketitle
\paragraph{Question and precise conclusion.}
The suggested critical four-dimensional Sobolev route would bound
the induced vertices of the same coefficient schur cumulant decomposition calculation by physical quadratic energy. Such an
argument first needs a COMPATIBLE coordinate/gauge representation.
The following actual relative-tree family proves that a universal
bound $|T_4(h)|\le C E(h)^2$ is false for raw tree coordinates.
It does not refute a bound on a completed gauge-invariant coefficient,
or on an induced quartic with a specifically proved exchange cancellation.

\paragraph{1. Original field and actual relative forest.}
Take the four-dimensional cube of side $R\ge6$, with all outer
boundary links fixed to the identity. In the spatial three-cube let
the Lie-algebra-valued function $\phi_0$ have only the values
\[
 \phi_0(3,2,2)=e_1,\quad \phi_0(3,3,2)=e_1+e_2,
 \quad\phi_0(2,3,2)=e_2,
\]
and vanish elsewhere. Here $e_1,e_2$ are two orthonormal SU(2) colors
in the quaternion normalization used throughout.
Set $\phi=R^{-1/2}\phi_0$ and define the physical linear field
$a_0(t,x)=\phi(x)$, $a_i=0$ for $i=1,2,3$.
It obeys the original outer boundary condition. Its plaquette energy is
\[
 E(a)=R\sum_{\text{spatial edges}}|\delta\phi|^2=20. \tag{1}
\]
The finite spatial gradient sum is $20$: vertex degree six gives $24$,
and the two nonzero neighboring inner products subtract $4$.

Use the genuine temporal relative forest: every interior column is
anchored at its lower boundary vertex, and all temporal links with
$t=0,\ldots,R-2$ are tree links. The last temporal link is a chord
to the fixed upper boundary. Set $\xi(t,x)=t\phi(x)$ at interior
vertices and $\xi=0$ on the outer boundary. The linear tree field
$h=a-d\xi$ is
\[
 h_0(t,x)=\begin{cases}0&t<R-1,\\R\phi(x)&t=R-1,\end{cases}
 \qquad h_i(t,x)=-t\,\delta_i\phi(x)\quad(0<t<R),             \tag{2}
\]
and $h_i=0$ on the temporal boundary faces. All tree components
vanish exactly, and $dh=da$, so $E(h)=20$. This is not an arbitrary
large pure-gauge field appended to an already fixed tree.

\paragraph{2. Actual Wilson quartic growth.}
On each spatial plaquette, $dh=0$. If its four oriented logarithms
are $X_1,\ldots,X_4$, the original product has logarithm
$\alpha^2\Omega_2+O(\alpha^3)$, where
$\Omega_2=\frac12\sum_{r<s}[X_r,X_s]$.
Its Wilson quartic is therefore $|\Omega_2|^2/2\ge0$.
On the indicated spatial square the logarithms are
$-t e_1/\sqrt R,-t e_2/\sqrt R,t e_1/\sqrt R,t e_2/\sqrt R$,
giving $2t^4/R^2$. The sum over $t=1,\ldots,R-1$ is at least
$cR^3$. The other spatial plaquettes cannot cancel it.

On ordinary temporal plaquettes, the two nonzero logarithms are
parallel and their difference is $\delta\phi$. Their total quartic
is bounded by $C/R$. The last temporal layer contains only finitely
many nonzero spatial stencils, each with logarithms $O(\sqrt R)$;
its absolute quartic is at most $CR^2$. Hence already
\[
 T_4(h)\ge cR^3-CR^2-C/R.                                  \tag{3}
\]
Every logarithm uses only colors one and two, so the original cubic
triple product is zero: $T_3(h)=0$.

An exact original-plaquette finite algebraic calculation gives the sharper formula
\[
 T_4(h)=R^3-\tfrac32R^2+\tfrac13R-\frac4{3R}.                \tag{4}
\]
This can be certified by finite polynomial arithmetic: before division
by $R^2$, the spatial part is $5\sum_{t=1}^{R-1}t^4$ and the temporal
part is a polynomial of degree at most four. Thus the total has degree
at most five. Exact evaluations at six distinct $R\ge6$ verify the
displayed polynomial. The finite calculation uses the complete original SU(2)
plaquette product, including the last layer; no continuum vertex or
infinite-volume identity is substituted.

\paragraph{3. Meaning for the proposed induced-vertex route.}
For the limiting Schur split with all coordinates retained, $H=1$
and the exchange term in the same coefficient schur cumulant decomposition calculation is absent. This actual example rules out
an unconditional raw-tree energy bound for that family of induced
representations. It also shows why a curl-controlled $L^4$ norm cannot
be assigned to arbitrary tree fields. At small $\alpha$ the field lies
in the original local chart; for instance $\alpha\sqrt R\to0$ suffices.
The obstruction is not caused by leaving the perturbative chart.

For a SPECIFIED nonempty block split, the negative exchange
$-\frac12\nabla_I T_3 C_I\nabla_I T_3$ may cancel part of this growth.
No failure for every such split is claimed. A Coulomb representation
can have critical Sobolev control, but moving from (2) to that
representation changes nonlinear coordinates and retained skeleton
sources. The source/frame/Jacobian contacts of the internal source tree to fp transport calculation and the two stage quotient and compensated covariance calculation and the exchange
must be carried through and estimated together. Neither the fact
that $T_3(h)=0$ nor the source-free shift of the relative background shift and its limit calculation removes those contacts.

Thus this lead does not presently give a valid quantitative bound
for the completed coarse remainder of the same coefficient schur cumulant decomposition calculation. The next necessary input is an
explicit compatible Coulomb/coarse representation with its compensated
vertices, or a bound on the original completed covariance of the two stage quotient and compensated covariance calculation itself.
The example is not a divergence of the full coefficient, and proves
nothing negative about existence, the nonlinear remainder, or S94.
\end{document}
